\thispagestyle{empty}
\section*{Abstract}
\label{sec:abstract}

With the rising popularity and easier access to 3D printers, more and more people are interested in designing and printing their own models. 
There are however still a few steps involved that require domain specific knowledge oftentimes creating a discrepancy between the wanted outcome and reality. 
Preparing a 3D model for printing requires the user to apply material presets and printing parameters after importing the model to a slicer software. 
Doing this once is a logical and necessary step, but when creating a 3D model from scratch or changing one, it might take several iterations until the result matches the desired outcome.
Thus the user has to repeatedly apply the aforementioned settings and therefore they have to remember them or write them down which creates friction in the workflow. 
Not only does this increase cognitive workload, leading to a higher frequency of errors and thus produces waste, it also often leads to frustration in beginner or casual makers. 
The proposed solution are so-called print hints added in the 3D modeling software itself, in the proof of concept realized using Fusion appearances. 
Upon opening a 3MF file in OrcaSlicer, the appearances are mapped to either material presets or modifiers. 
By giving the user the option to map manually while also saving and loading previous mappings, the solution provides an automation of the criticized workflow while keeping the freedom to easily change an entire group of objects in one simple step. 

\section*{Kurzdarstellung}
\label{sec:kurzdarstellung}

\todo translation
